\section{Task assignment}
\label{sec:assignment}

\begin{lede}
Which robot gets which job. A greedy match over a cost that knows about
distance, congestion, one-way aisles and how long the job has queued.
\end{lede}

\subsection{Directional delay}
\label{sec:assignment-delay}

\begin{equation}
  T_{\mathrm{dir}}(s, g) = \sum_{a \in \mathcal{A}(\mathrm{route}(s, g))} \mathrm{delay}(a),
\end{equation}
\begin{equation}
  \mathrm{delay}(a) =
  \begin{cases}
    0 & a \text{ is \textsc{open}},\\
    \mathrm{occ}(a) + 1 & a \text{ is \textsc{draining}},\\
    \tfrac{1}{2}T_{\min} + \mathrm{occ}(a) & \text{direction mismatches},\\
    0 & \text{direction already matches.}
  \end{cases}
\end{equation}
The third case is the cost of waiting for a drain and a flip. Computed only
when \param{direction\_control} is \code{"aisle"}, with each aisle counted once
even if the route touches it repeatedly.

\subsection{Blocking estimate}

\begin{equation}
  B(i, \tau) = \frac{|\{\text{bottleneck vertices on } \mathrm{route}(x_i, p_\tau)\}|}
    {\max(1, |\mathrm{route}|)}.
\end{equation}

\subsection{The assignment cost}
\label{sec:assignment-cost}

\begin{equation}
  J(i, \tau) = a\, d(x_i, p) + b\, d(p, d) + g\, C + \delta\, W
    + e\, T_{\mathrm{dir}} + z\, B.
  \label{eq:assignment}
\end{equation}

\begin{center}\small
\begin{tabular}{@{}lll@{}}
\toprule
\textbf{Term} & \textbf{Meaning} & \textbf{Coefficient} \\
\midrule
$a\, d(x_i, p)$ & robot to pickup distance & \param{assign\_alpha\_to\_pickup} $= 1.0$ \\
$b\, d(p, d)$ & pickup to delivery distance & \param{assign\_beta\_pickup\_to\_delivery} $= 0.5$ \\
$g\, C$ & route congestion, if congestion-aware & \param{assign\_gamma\_congestion} $= 12.0$ \\
$\delta\, W$ & task waiting time, capped at 60 & \param{assign\_delta\_waiting} $= -0.5$ \\
$e\, T_{\mathrm{dir}}$ & directional delay, \Cref{sec:assignment-delay} & \param{assign\_eta\_direction} $= 1.0$ \\
$z\, B$ & bottleneck blocking, if congestion-aware & \param{assign\_zeta\_blocking} $= 1.0$ \\
\bottomrule
\end{tabular}
\end{center}

Note that $\delta = -0.5$ is \textbf{negative}: a longer-waiting task
\emph{lowers} its own cost and is preferred over a fresher, possibly nearer
one. This is the assignment layer's own fairness mechanism, separate from and
additive to the priority function's waiting term (\Cref{sec:priority}). The
cost returns $\INF$ if either leg is unreachable, so an infeasible pairing is
never chosen.

\begin{caveat}
Two coefficients here were recalibrated rather than redesigned, and both
mattered. \param{assign\_gamma\_congestion} was $2.0$ against distances of
10--40, so the entire congestion term was worth about four cost units and could
not influence a match --- a congestion-aware assignment that cannot change an
assignment is not a mechanism. And \param{assign\_waiting\_cap} bounds a term
that otherwise grows without limit in a lifelong run: uncapped, waiting
dominates distance and congestion within about a hundred steps and the match
degenerates to oldest-task-first, which is exactly what hypothesis H4 was
reporting before the cap existed.
\end{caveat}

% worked-example: assignment
\begin{workedexample}
\textbf{Worked example.} Same run, timestep 18. One unassigned task: pickup (0,0), delivery (15,0), released at t = 6 so it has been waiting 12 steps. Two candidate robots, costed by \code{TaskAssigner.assignment\_cost}:

{\fontsize{9.0}{10.8}\selectfont
\begin{verbatim}
robot           a·d(r,p)  b·d(p,d)  g·C    d·W  e·T_dir  z·B  J(i,τ)
--------------  --------  --------  -----  ---  -------  ---  ------
robot 16 (0,4)  4         7.5       2.4    −6   0        0    7.9
robot 18 (0,8)  8         7.5       2.667  −6   0        0    12.17
\end{verbatim}
}

Robot 16 wins at J = 7.9. Two columns are worth pausing on. \code{d(pickup, delivery)} is identical for both robots, because it does not depend on which robot goes — it is in the cost only so that a queue of tasks is ordered by total work rather than by how near its pickup happens to be. And \code{δ} = −0.5 is \emph{negative}, so the waiting term is a discount: the longer a task has been queued the cheaper it looks, which is how ageing gets into a greedy match.

That discount is capped at \code{assign\_waiting\_cap} = 60. Uncapped, \code{waiting\_time} grows without bound in a lifelong run and dwarfs both distance and congestion within about a hundred steps; the match degenerates into oldest-task-first, and a congestion-aware assignment can no longer show any effect at all — which is exactly what H4 reported before the cap existed (\Cref{sec:assignment-cost}).
\end{workedexample}
% /worked-example

\subsection{Greedy matching}

\code{assign\_tasks\_greedily} first pre-filters the unassigned pool to the
oldest
\[
  \max\bigl(\param{assignment\_candidate\_limit},\ |\text{free robots}|\bigr)
\]
tasks, sorted by $(t_{\mathrm{rel}}, \mathrm{id})$ --- a cheap bound, so cost is
never computed against an unbounded backlog. It then repeatedly takes the
cheapest feasible pair:

\begin{algorithm}[H]
\caption{\code{assign\_tasks\_greedily}}
\begin{algorithmic}[1]
\While{free robots and candidate tasks remain}
  \State $(i^\ast, \tau^\ast) \gets \arg\min_{(i, \tau) \text{ feasible}} J(i, \tau, t)$
  \State assign $i^\ast \leftarrow \tau^\ast$; remove both from their pools
\EndWhile
\end{algorithmic}
\end{algorithm}

This is greedy minimum-cost matching, $O(RT)$ per round for $R$ robots and $T$
candidate tasks --- not the Hungarian algorithm, so it is locally greedy at
each pairing rather than globally optimal. It is held \emph{identical} across
every planner compared in \Cref{sec:results}, including both external
baselines, so that comparisons isolate the movement layer.
